\documentclass[11pt]{article}
\usepackage[T1]{fontenc}
\usepackage[utf8]{inputenc}
\usepackage{lmodern}
\usepackage[margin=1in]{geometry}
\usepackage{amsmath,amssymb,amsthm,mathtools}
\usepackage{microtype}
\usepackage[numbers,sort&compress]{natbib}
\usepackage[hidelinks]{hyperref}
\usepackage{url}
\hypersetup{
  pdftitle={Interval shape and endpoint collisions for totally positive shifted even powers},
  pdfauthor={Xiaojian Zeng},
  pdfsubject={The interval-shape assertion of Vishnyakova's Conjecture 5.8},
  pdfkeywords={total positivity, Euler-Frobenius polynomials, real zeros, Bernstein certificates}}
\newtheorem{theorem}{Theorem}[section]
\newtheorem{lemma}[theorem]{Lemma}
\newtheorem{proposition}[theorem]{Proposition}
\newtheorem{corollary}[theorem]{Corollary}
\theoremstyle{remark}
\newtheorem{remark}[theorem]{Remark}
\newcommand{\R}{\mathbb{R}}
\newcommand{\Q}{\mathbb{Q}}
\newcommand{\T}{\mathcal{T}}
\newcommand{\Lop}{\mathcal{L}}
\DeclareMathOperator{\sgn}{sgn}
\DeclareMathOperator{\disc}{disc}
\title{Interval shape and endpoint collisions for\protect\\
totally positive shifted even powers}
\author{Xiaojian Zeng\thanks{PBC School of Finance, Tsinghua University,
100190. Email: \href{mailto:xiaojianzeng@foxmail.com}{xiaojianzeng@foxmail.com}.}}
\date{}
\begin{document}
\maketitle
\begin{abstract}
For an even integer $n\ge2$, let $M_n$ be the set of real shifts $t$ for
which the sequence $((k+t)^n)_{k\ge0}$ has a totally nonnegative Toeplitz
matrix. We prove that $M_n=[-c_n,1+c_n]$, where
$(\sqrt2-1)/2\le c_n<1-1/\sqrt2$. Each exterior endpoint corresponds to
one double zero of the finite-difference polynomial, with all other zeros
simple and with nondegenerate splitting toward the admissible interval.
The proof uses a coefficient-constrained two-step differential operator
and a uniform sign estimate based on Cauchy--Schwarz. It also locates all
zeros of every odd-degree shifted-power transform for $-3/10\le t<0$.
This proves the interval-shape assertion of Vishnyakova's Conjecture 5.8.
The successive-nesting assertion $c_n\le c_{n+2}$ remains unresolved here,
so the full conjecture is not settled.
\end{abstract}

\section{Statement and context}

For a real sequence $(a_j)_{j\ge0}$, extend $a_j=0$ for $j<0$.
Following the terminology of Vishnyakova \cite{Vishnyakova2026}, call the
sequence totally positive if every finite minor of
$(a_{j-i})_{i,j\ge0}$ is nonnegative. Thus zero minors are allowed.
For an integer $n\ge1$ put
\[
 M_n=\{t\in\R: ((k+t)^n)_{k\ge0}\text{ is totally positive}\}.
\]
Vishnyakova \cite[Example 5.7 and Conjecture 5.8]{Vishnyakova2026}
asks whether, for every $k\ge1$,
\begin{equation}\label{eq:original}
 M_{2k}=[-c_{2k},1+c_{2k}],\qquad 0<c_{2k}<2,
 \qquad c_{2k}\le c_{2k+2}.
\end{equation}
We establish the interval and bound assertions, together with a description
of the boundary roots. The last assertion of \eqref{eq:original} remains
unresolved in this paper.

Define $\Lop\binom{x}{j}=z^j$ and write
\begin{equation}\label{eq:Q}
 Q_{n,t}(z)=\Lop((x+t)^n)
 =\sum_{j=0}^{n}z^j\sum_{\ell=0}^{j}
       (-1)^\ell\binom{j}{\ell}(j-\ell+t)^n.
\end{equation}
This polynomial has degree $n$ and leading coefficient $n!$.
The generating-function characterization of total positivity
\cite{Edrei1953}, applied as in \cite[Section 1]{Vishnyakova2026}, gives
\begin{equation}\label{eq:criterion}
 t\in M_n\quad\Longleftrightarrow\quad
 \text{all $n$ zeros of $Q_{n,t}$, with multiplicity, lie in $[-1,0]$}.
\end{equation}
Indeed,
\begin{equation}\label{eq:gen}
 \sum_{k\ge0}(k+t)^n x^k
 =\frac{E_{n,t}(x)}{(1-x)^{n+1}},\qquad
 E_{n,t}(x)=(1-x)^nQ_{n,t}\!\left(\frac{x}{1-x}\right).
\end{equation}
If $Q_{n,t}(z)=n!\prod_{i=1}^n(z+a_i)$ with $0\le a_i\le1$, then
$E_{n,t}(x)=n!\prod_i(a_i+(1-a_i)x)$, so its coefficients are nonnegative
and its finite zeros nonpositive. The converse follows from the same
generating-function criterion. A factor with $a_i=1$ is constant in the
numerator; the possible degree drop loses no endpoint case. Also
$E_{n,t}(1)=n!$, so the denominator in \eqref{eq:gen} has no cancellation.

\begin{theorem}\label{thm:main}
For every even integer $n\ge2$ there exists $c_n$ such that
\begin{equation}\label{eq:main}
 M_n=[-c_n,1+c_n],\qquad
 \frac{\sqrt2-1}{2}\le c_n<1-\frac1{\sqrt2}.
\end{equation}
At $t=-c_n$, $Q_{n,t}$ has exactly one double zero $r\in(-1,0)$.
All its other zeros are simple and smaller than $r$, and
\begin{equation}\label{eq:collision-signs}
 \partial_z^2Q_{n,t}(r)>0,\qquad \partial_tQ_{n,t}(r)<0.
\end{equation}
The endpoint $t=1+c_n$ has the reflected configuration. At either
endpoint the double zero splits into two simple real zeros on the inward
side and into a nonreal conjugate pair on the outward side.
\end{theorem}

The numerator $E_{n,t}$ in \eqref{eq:gen} is the Euler--Frobenius
polynomial $P_{n,\rho}$ with $\rho=t$ in Janson's notation
\cite[Section 1]{Janson2013}. The standard simple-negative-root assertion in
\cite[Remark A.2]{Janson2013} concerns $0<\rho\le1$.
The result below supplies an exterior parameter strip. A related
interlacing lemma of Karp and Vishnyakova
\cite[Lemma 3]{KarpVishnyakova2025} assumes that every input numerator
zero avoids $[0,1]$. Our positive $Q$-zero $\lambda$ corresponds to the
numerator zero $\lambda/(1+\lambda)\in(0,1)$, so direct iteration of that
lemma does not give our induction.

Dubeau and Savoie \cite[Theorems 12--13]{DubeauSavoie1998} study
two-degree root comparisons for their recurrence (G), whose differential
operator is $(x^2-1)D+2r_nx$ with $r_n>0$. Under $x=2z+1$, twice our
operator becomes $(x^2-1)D+x+(2t-1)$. Its constant term persists on the
exterior strip. These particular hypotheses therefore do not directly
supply the theorem proved here. The contribution is an all-degree
exterior induction and the exclusion of separated admissible components.
No equivalent was found in these inspected statements; this comparison
does not exhaust the older literature or establish priority.

\section{A coefficient-constrained two-step operator}

For fixed $t$ define
\begin{equation}\label{eq:T}
 \T_t p=z(1+z)p'+(z+t)p.
\end{equation}

\begin{theorem}\label{thm:operator}
Let $-3/10\le t<0$, and let $h$ have even degree $d\ge2$, positive
leading coefficient, and $d$ simple zeros in $(-1,0)$. Put
\[
 S=\frac{h'(0)}{h(0)},\qquad
 \delta=\frac{1+2t}{(1+t)^2},\qquad
 \Lambda=\frac1{S+\delta}.
\]
For every $0<\lambda\le\Lambda$, the polynomial
$q=\T_t^2[(z-\lambda)h]$ has $d+2$ simple zeros in $(-1,0)$ and one
simple positive zero. These are all its zeros.
\end{theorem}

The coefficient constraint includes equality. It can be read directly
from the output: for $p=(z-\lambda)h$,
\begin{equation}\label{eq:coefficient}
 [z]\T_t^2p=(1+t)^2h(0)\{1-\lambda(S+\delta)\}.
\end{equation}
To check this, if $p=p_0+p_1z+\cdots$, then
$[z]\T_t^2p=(1+t)^2p_1+(1+2t)p_0$.
Thus the hypothesis is precisely nonnegativity of this linear coefficient.

\begin{proof}[Proof of Theorem \ref{thm:operator}]
Write $h(z)=c\prod_{i=1}^d(z+a_i)$, where $c>0$ and $0<a_i<1$.
Set
\begin{gather*}
 y=\Lambda/2,\qquad r_i=y/a_i,\qquad
 R=\sum_i r_i=yS=\tfrac12-\delta y,\qquad
 \eta=\frac{y}{y+\lambda},\\
 w_i=\frac{r_i}{1-r_i},\qquad W=\sum_i w_i,\qquad V=\sum_i w_i^2.
\end{gather*}
Since $S=\sum_i a_i^{-1}>d\ge2$, we have
\begin{equation}\label{eq:bounds}
 \frac{40}{49}\le\delta\le1,\qquad
 0<y<Y:=\frac1{2(2+\delta)}<\frac5{28},\qquad
 0<R<\frac12,\qquad \eta\ge\frac13.
\end{equation}
Also $y<1/(2S)<a_i/2$, so $y<r_i<1/2$.
In particular, $-y$ lies to the right of every negative zero of $p$, and
$p(-y)<0$. We will prove $q(-y)>0$.

First,
\begin{equation}\label{eq:Wlower}
 W\ge\frac{R}{1-y}.
\end{equation}
Weighted Cauchy--Schwarz gives the further estimate
\[
 W^2=\left(\sum_i\frac{r_i}{1-r_i}\right)^2
 \le\left(\sum_i r_i\right)
      \left(\sum_i\frac{r_i}{(1-r_i)^2}\right)=R(W+V),
\]
and hence
\begin{equation}\label{eq:cauchy}
 W^2-V\le W-\frac{1-R}{R}W^2.
\end{equation}
With $a=z(1+z)$, direct differentiation in \eqref{eq:T} yields
\[
 \T_t^2p=a^2p''+a(1+4z+2t)p'
              +\{2z^2+(1+2t)z+t^2\}p.
\]
At $z=-y$, logarithmic differentiation gives
\[
 \frac{p'}p=\frac{W-\eta}{y},\qquad
 \frac{p''}p=\frac{W^2-V-2\eta W}{y^2}.
\]
For $B=1-4y+2t$ and $C=2y^2-(1+2t)y+t^2$, this becomes
\begin{equation}\label{eq:quotient}
 \frac{q(-y)}{p(-y)}
 =(1-y)^2(W^2-V-2\eta W)-(1-y)B(W-\eta)+C.
\end{equation}
The coefficient of $\eta$ is negative, because \eqref{eq:Wlower} implies
\[
 B-2(1-y)W\le B-2R=2t+2(\delta-2)y<0.
\]
Consequently, replacing $\eta$ by $1/3$ gives an upper bound even when
$\lambda=\Lambda$. Let
\[
 K=(1-y)\left(\frac53-\frac{14}3y+2t\right)>0.
\]
The bracket is greater than
$5/3-(14/3)(5/28)-3/5=7/30$. By \eqref{eq:cauchy},
\[
 \frac{q(-y)}{p(-y)}\le F(W)+\frac{(1-y)B}{3}+C,
 \quad F(W)=(1-y)^2\left(W-\frac{1-R}{R}W^2\right)-KW.
\]
For $W\ge R$,
$F'(W)\le(1-y)^2(2R-1)-K<0$. We may therefore substitute the
lower bound $W=R/(1-y)$. After using $R=1/2-\delta y$, we obtain
\begin{align}
 \frac{q(-y)}{p(-y)}&\le U_t(y),\label{eq:Ubound}\\
 U_t(y)&=t^2-\frac t3-\frac14
 +\left(2\delta t+\frac{2\delta}3-\frac{8t}3-\frac56\right)y
 +\left(\delta^2-\frac{11\delta}3+\frac{10}3\right)y^2.
 \label{eq:U}
\end{align}
The quadratic coefficient is $(\delta-5/3)(\delta-2)>0$.
It suffices by convexity to check $U_t(0)<0$ and $U_t(Y)<0$.
The first value is at most $-3/50$. For the second,
\begin{equation}\label{eq:N}
 U_t(Y)=\frac{N(t)}{3(2t^2+6t+3)^2},\qquad
 N(t)=12t^6+60t^5+89t^4+12t^3-60t^2-40t-7.
\end{equation}
Its denominator is positive since
$2t^2+6t+3=(1+t)^2(2+\delta)>0$.

Here a short exact Bernstein identity supplies the sign on the entire
interval. Set $b_{j,6}(u)=\binom6j u^j(1-u)^{6-j}$. Then
\begin{equation}\label{eq:bernsteinN}
 N\!\left(\frac3{10}(u-1)\right)=\sum_{j=0}^6 B_j b_{j,6}(u),
 \qquad 0\le u\le1,
\end{equation}
where
\[
 (B_0,\ldots,B_6)=\left(
 -\frac{17519}{125000},-\frac{273}{500},-\frac{58837}{50000},
 -\frac{10481}{5000},-\frac{84}{25},-5,-7\right).
\]
These identities can be checked by expansion. The basis functions are
nonnegative and sum to one, so $N(t)<0$ on $[-3/10,0]$.
Thus \eqref{eq:Ubound} is strictly negative and
\begin{equation}\label{eq:extra-sign}
 q(-y)>0.
\end{equation}

It remains to count the roots. Order the negative zeros of $p$ as
$\beta_1<\cdots<\beta_d$ and put $g=\T_t p$. Off its poles,
\begin{equation}\label{eq:H}
 \frac{g(z)}{z(1+z)p(z)}=H(z)
 =\sum_{i=1}^d\frac1{z-\beta_i}+\frac1{z-\lambda}
       +\frac t z+\frac{1-t}{1+z}.
\end{equation}
On $(-1,\beta_d)$, away from the poles, $H'(z)<0$.
Indeed its only positive summand is $-t/z^2$, and
$1/(z-\beta_d)^2>1/z^2>-t/z^2$.
The pole limits show that each interval
\[
 (-1,\beta_1),\ (\beta_1,\beta_2),\ldots,
 (\beta_{d-1},\beta_d)
\]
contains exactly one simple zero $\xi_i$ of $g$.
At this zero,
$g'(\xi_i)=\xi_i(1+\xi_i)p(\xi_i)H'(\xi_i)$, so
\[
 \sgn g'(\xi_i)=\sgn p(\xi_i)=(-1)^i,
 \qquad \sgn q(\xi_i)=(-1)^{i+1}.
\]
Moreover,
$q(-1)=(t-1)^2p(-1)<0$, $q(0)=t^2p(0)<0$, and
$\xi_d<\beta_d<-y<0$. These signs and \eqref{eq:extra-sign} give
one zero of $q$ in each of
\[
 (-1,\xi_1),\ (\xi_1,\xi_2),\ldots,(\xi_{d-1},\xi_d),\quad
 (\xi_d,-y),\quad(-y,0),\quad(0,\infty).
\]
The last interval uses the positive leading coefficient of $q$.
Since $\deg q=d+3$, the exhibited $d+3$ distinct zeros exhaust its degree
and are all simple.
\end{proof}

The sign argument uses the elementary convex-hull property of the
Bernstein basis. Polynomial arithmetic in this basis is developed by
Farouki and Rajan \cite{FaroukiRajan1988}. A related exact-arithmetic
application appears in Zeng \cite{ZengLens2026}: the proof of Theorem 9,
equation (8.18), certifies a degree-six parameter polynomial by seven
positive rational Bernstein coefficients in a spectral inequality.
Here both coefficient lists are displayed in full, so the sign certificates
require only rational polynomial expansion.

\section{Odd shifted powers outside the classical interval}

The identity
$x\binom{x}{j}=j\binom{x}{j}+(j+1)\binom{x}{j+1}$ gives
\begin{equation}\label{eq:recurrence}
 Q_{0,t}=1,\qquad Q_{n+1,t}=\T_tQ_{n,t},\qquad
 \partial_tQ_{n,t}=nQ_{n-1,t}.
\end{equation}
The last equality follows directly by differentiating before applying
$\Lop$. The recurrence and finite differences also give
\begin{gather}
 Q_{n,t}(0)=t^n,\quad Q_{n,t}(-1)=(t-1)^n,\quad
 [z^n]Q_{n,t}=n!,\quad [z]Q_{n,t}=(1+t)^n-t^n,
 \label{eq:values}\\
 Q_{n,t}(-1-z)=(-1)^nQ_{n,1-t}(z).
 \label{eq:reflection}
\end{gather}
The reflection follows by substituting into \eqref{eq:recurrence}, starting
with $Q_{0,t}=1$.

\begin{theorem}\label{thm:odd}
For every odd integer $m\ge1$ and every $-3/10\le t<0$, the polynomial
$Q_{m,t}$ has $m-1$ simple zeros in $(-1,0)$ and one simple positive zero.
\end{theorem}

\begin{proof}
The case $m=1$ is $Q_{1,t}=z+t$. For $m=3$,
\begin{align*}
 Q_{3,t}&=6z^3+6(1+t)z^2+(3t^2+3t+1)z+t^3,\\
 D_3(t):=\disc_zQ_{3,t}
 &=-12(18t^6-54t^5+36t^4+18t^3-12t^2-6t-1).
\end{align*}
The identity
\begin{equation}\label{eq:bernsteinD}
 D_3\!\left(\frac3{10}(u-1)\right)=\sum_{j=0}^6 C_jb_{j,6}(u)
\end{equation}
has coefficient vector
\[
 (C_0,\ldots,C_6)=\left(
 \frac{495087}{125000},\frac{103179}{25000},\frac{11616}{3125},
 \frac{10209}{2500},\frac{708}{125},\frac{42}{5},12\right).
\]
Hence $D_3(t)>0$ on $[-3/10,0]$, so the cubic has three distinct real
zeros. For $t<0$ its nonconstant coefficients are positive and its
constant coefficient negative. Descartes' rule gives exactly one positive
zero. Reflection \eqref{eq:reflection} excludes a zero at or below $-1$,
since every coefficient of $Q_{3,1-t}$ is positive. The other two zeros
therefore lie in $(-1,0)$.

Suppose the result holds for odd $m\ge3$, and write
$Q_{m,t}=(z-\lambda)h$ with $\lambda>0$. Then $h$ satisfies the hypotheses
of Theorem~\ref{thm:operator}. Its two-step output is $Q_{m+2,t}$, whose
linear coefficient is
\[
 (1+t)^{m+2}-t^{m+2}>0.
\]
Equation \eqref{eq:coefficient}, with $h(0)>0$, implies
$\lambda<1/(S+\delta)$. Theorem~\ref{thm:operator} completes the induction.
\end{proof}

The positivity used to obtain the coefficient constraint is an explicit
coefficient identity, independent of the output root configuration.
No assertion that odd shifted powers in this exterior strip are totally
positive is involved: their transform has a positive zero.

\section{The even-degree admissible interval}

Write
\[
 b=1-1/\sqrt2,\qquad c_2=(\sqrt2-1)/2.
\]
We first record the elementary restrictions and inclusions.

\begin{lemma}\label{lem:bounds}
For every even $n\ge2$, $M_n$ is compact and
\[
 [-c_2,1+c_2]\subseteq M_n\subset(-b,1+b).
\]
For $0<t<1$, all zeros of $Q_{n,t}$ are simple and lie in $(-1,0)$.
\end{lemma}

\begin{proof}
If $t\in M_n$, \eqref{eq:criterion} implies
$[z]Q_{n,t}=(1+t)^n-t^n\ge0$, so $t\ge-1/2$.
For $-1/2\le t<0$, the Toeplitz minor with rows $0,1$ and columns $1,2$
requires
\[
 (1+t)^{2n}\ge t^n(2+t)^n.
\]
As $n$ is even, this is equivalent to
$(1+t)^2\ge-t(t+2)$, whence $t\ge-b$.
At $t=-b$, let $a_j=(j-b)^n$. Then $a_1^2=a_0a_2>0$, and the
minor with rows $0,1,2$ and columns $1,2,3$ is
\begin{align*}
 \det\begin{pmatrix}a_1&a_2&a_3\\a_0&a_1&a_2\\0&a_0&a_1\end{pmatrix}
 &=a_1^3-2a_0a_1a_2+a_0^2a_3\\
 &=\frac{a_0^2}{a_1}(a_1a_3-a_2^2)<0.
\end{align*}
The strict sign follows from
$(1-b)(3-b)-(2-b)^2=-1$. Thus $-b\notin M_n$.
Reflection gives the upper restriction. Each defining minor is continuous
in $t$, so $M_n$ is closed and bounded, hence compact.

For the lower inclusion,
\[
 Q_{2,t}=2z^2+(1+2t)z+t^2,\qquad
 \disc_zQ_{2,t}=1+4t-4t^2.
\]
On the discriminant interval the two roots are nonpositive; reflection
makes them at least $-1$. Thus
$M_2=[-c_2,1+c_2]$, as in \cite[Example 5.3]{Vishnyakova2026}.
Wagner's multiplication theorem \cite{Wagner1992}, also stated as
Theorem C in \cite{Vishnyakova2026}, says that the class of polynomials
interpolating totally positive sequences is closed under polynomial
multiplication. Applying it repeatedly to $(x+t)^2$ yields $M_2\subseteq M_n$.

Finally, if $0<t<1$ and a polynomial $p$ has simple zeros in $(-1,0)$,
then
\[
 \frac{\T_tp}{z(1+z)p}=\frac{p'}p+\frac t z+\frac{1-t}{1+z}
\]
is strictly decreasing between its poles. All residues are positive,
and the pole limits give one simple zero in every interval between
successive poles from $-1$ to $0$. Induction from $Q_{1,t}=z+t$ proves
the assertion. Closedness includes $t=0,1$ in $M_n$.
\end{proof}

\begin{lemma}\label{lem:collision}
Let $n\ge2$ be even, let $-b<t_0<0$, and suppose $Q_{n,t_0}$ is
admissible and has a multiple zero. It has exactly the root configuration
and derivative signs in \eqref{eq:collision-signs}. Admissibility holds
immediately to the right of $t_0$ and fails immediately to its left.
\end{lemma}

\begin{proof}
For $n\ge4$, put $f=Q_{n,t_0}$ and $g=Q_{n-1,t_0}$.
Since $b<3/10$, Theorem~\ref{thm:odd} gives $n-2$ simple negative zeros
$\beta_1<\cdots<\beta_{n-2}$ of $g$ and one simple positive zero
$\lambda$. The strictly decreasing expression \eqref{eq:H}, now for
$f=\T_{t_0}g$, gives exactly one simple zero of $f$ in each of
\[
 (-1,\beta_1),\ (\beta_1,\beta_2),\ldots,
 (\beta_{n-3},\beta_{n-2}).
\]
No $\beta_i$ is a zero of $f$, since
$f(\beta_i)=\beta_i(1+\beta_i)g'(\beta_i)\ne0$.
Neither endpoint $-1,0$ is a zero. Admissibility and the degree count
therefore force every multiple zero into $(\beta_{n-2},0)$, where only
two degrees remain. There is exactly one double zero $r$, and all other
zeros $\gamma$ of $f$ are simple and smaller than $r$. Hence
\[
 f''(r)=2n!\prod_{\gamma}(r-\gamma)>0,\qquad
 \partial_tQ_{n,t_0}(r)=ng(r)<0.
\]
For the second sign, all negative-root factors of $g$ are positive at
$r$, and its positive-root factor is negative.

For $n=2$ the only negative multiple-root parameter is $t_0=-c_2$.
Its double root is $r=-(1+2t_0)/4\in(-1,0)$, and direct differentiation
gives $\partial_z^2Q_{2,t_0}(r)=4$ and
$\partial_tQ_{2,t_0}(r)=t_0-1/2<0$.

To determine the direction of the collision, put $s=t-t_0$ and factor
off the persistent simple roots near $s=0$. The remaining monic factor
has the form
\[
 (z-r)^2+A(s)(z-r)+B(s),\qquad A(0)=B(0)=0.
\]
The implicit function theorem for the simple roots makes $A,B$ real
analytic, and differentiation at $s=0$ gives
\[
 B'(0)=\frac{2\partial_tQ_{n,t_0}(r)}{\partial_z^2Q_{n,t_0}(r)}<0.
\]
Thus $(A^2-4B)'(0)=-4B'(0)>0$. The two roots are distinct and real for
small $s>0$, and are nonreal conjugates for small $s<0$.
All persistent simple roots, and the split real roots, stay in $(-1,0)$.
This proves both the nondegeneracy and the asserted orientation.
\end{proof}

\begin{proof}[Proof of Theorem \ref{thm:main}]
The discriminant of $Q_{n,t}$ is a polynomial in $t$ that is nonzero at
$t=1/2$ by Lemma~\ref{lem:bounds}. It has finitely many zeros.
Away from those zeros in $(-b,0)$, the roots vary continuously, remain
simple, and cannot cross $-1$ or $0$ by \eqref{eq:values}. Thus
admissibility is locally constant there, and there are only finitely
many admissible components in $(-b,0]$.

Every admissible parameter with simple roots is interior. By
Lemma~\ref{lem:collision}, every negative admissible collision has
admissibility to its right and failure to its left. It cannot be a
right boundary or an isolated admissible point. The lower inclusion
in Lemma~\ref{lem:bounds} supplies a component meeting $0$. Any further
component would have a negative right boundary, necessarily admissible
by closedness, contradicting the preceding alternatives.
Consequently $M_n\cap(-b,0]=[-c_n,0]$. Compactness and the exclusion of
$-b$ give $c_n<b$, while $M_2\subseteq M_n$ gives $c_n\ge c_2$.
Reflection and $[0,1]\subseteq M_n$ prove \eqref{eq:main}.

At $-c_n$ the roots cannot all be simple, because that would make the
parameter interior. Lemma~\ref{lem:collision} gives the endpoint
configuration and signs. Equation \eqref{eq:reflection} supplies the
configuration at $1+c_n$ and reverses the parameter direction.
\end{proof}

\section{Exact checks, assistance, and remaining scope}

The proof above is deductive. Its two polynomial sign certificates
\eqref{eq:bernsteinN} and \eqref{eq:bernsteinD} are finite identities
valid throughout the stated parameter interval, not samples of it.
They can be verified by expansion over $\Q$ or by matching coefficients.
Equivalently, for $P(u)=\sum_{k=0}^6p_ku^k$, its degree-six Bernstein
coefficients are
\[
 B_j=\sum_{k=0}^j p_k\frac{\binom jk}{\binom6k},\qquad 0\le j\le6.
\]
This formula follows by expanding $(u+(1-u))^{6-k}$ after multiplying by
$u^k$.

The accompanying verification programs use Python 3.10 and SymPy 1.14
\cite{Meurer2017}. They check ten symbolic identities and 72 rational
operator inputs, including 36 inputs with $\lambda=\Lambda$. A second
implementation reconstructs the two Bernstein lists by rational
interpolation, checks those 72 inputs, and checks 27 further inputs and
six odd shifted powers. The latter have degrees $15,17,21$, each at
$t=-3/10$ and $t=-1/1000$. The operator inputs, exact coefficients,
evaluation signs, and root counts are recorded with the programs.
Rational Euclidean Sturm chains and exact polynomial gcds check location
and simplicity; no floating-point root decision enters these checks.
The finite tests check algebra and implementation and do not replace
the all-degree argument.

\paragraph{Assistance and reproducibility.}
OpenAI language-model assistance was used in proof development,
counterexample exploration, code generation, literature searches,
mathematical checking, and manuscript preparation. A separate model
context reconstructed the argument and recomputed the exact certificates
with a different implementation. This was an internal check, not external
peer review or proof-assistant formalization. The accompanying source
package contains the exact verifier inputs and a reproduction guide;
the two sign certificates needed in the proof are also displayed above.

Theorem~\ref{thm:main} concerns every even degree, but it compares no two
successive boundary constants. The inequality $c_n\le c_{n+2}$ in
\eqref{eq:original} remains unproved here, as does the contextual question
of a possible limit. The endpoint $-3/10$ in
Theorem~\ref{thm:operator} is a sufficient endpoint; no sharpness is
asserted. The quadratic inclusion from multiplication closure alone does
not imply successive nesting.

\bibliographystyle{plainnat}
\small
\bibliography{references}

\end{document}
